\section{Constrained affine sampling}
\label{sampling:constrained-section}

We now require sampled affine boxes to carry small real lifts of
specified polynomials, with integral polynomial complements.
An independent reference law on coefficient arrays describes the
desired lifts. Evaluating its density on the coefficients of
$C_h(x+Vt)$ then tilts the spatial law of $(x,V)$. These are distinct
laws: the actual polynomial coefficients are determined by the
sampled affine map, not sampled independently of it.

The construction will support two later arguments. Cube comparison
in Section~\ref{cube:section} transfers path-dependent detection and
gives finite approximation with pre bounds. Positive score recovery
in Section~\ref{cube:transfer-section} uses a nonnegative comparison
weight for the same ensemble; only its final patch may incur the
specified late costs. We state both outputs before constructing the
charts and coefficient tilt, and prove here the normalization and
residue estimates that those two arguments share.

\subsection{The constrained sampler and its parameter order}

Let $U=\prod_{i=1}^n\{0,\ldots,H_i-1\}$, with $n\ge1$.  For $1\le h\le s$, let
$W_h\subseteq\R^{d_h}$ be a rational subspace and let
$C_h:\Z^n\to W_h$ be an ordinary real polynomial of degree at most $h$.
Neither its coefficients nor the height of $W_h$ is bounded at this stage.
Write $\Lambda_h=W_h\cap\Z^{d_h}$ and fix $c_h\in W_h/\Lambda_h$.
Choose a representative in $W_h$, also denoted $c_h$.
Within the chart used below, write
\begin{equation}\label{sampling:lift}
 C_h(u)-c_h=y_h(u)+\beta_h(u),\qquad
 \beta_h(u)\in\Z^{d_h},\qquad |y_h(u)|_\infty<\tfrac14.
\end{equation}
The chart makes this decomposition unique whenever it exists.

\begin{definition}[Rank at scale $H$]\label{sampling:rank}
The system has rank at least $R$ if, for every $h$ and every integer row
$a\in\Z^{d_h}$ with $\|a\|\le R$ and $a|_{W_h}\ne0$, the homogeneous
degree-$h$ part of $aC_h$ cannot be written as a polynomial with coefficient
bounds $R/H^\nu$ plus a rational polynomial whose coefficients have a common
denominator at most $R$.
\end{definition}

The sampler has three stages of parameters.
First fix the degree and structural dimensions $n,\sum_hd_h$.
The sampling dimension $m$ and every block count are fixed polynomials
of these dimensions, independent of testing accuracy, recursive output
budgets, and chart widths.

Next choose a budget $b_{\mathrm{pre}}\ge2$ for the fixed dimensions,
testing complexities, logarithms of inverse chart widths, and requested
logarithmic inverse accuracies.  A \emph{pre bound} means a bound of the
following size, with $C$ depending only on the fixed degrees:
\[
\begin{array}{c|c}
\text{quantity} & \text{allowed bound}\\ \hline
\text{counts, caps, and coefficient masses}
  & \exp((2+b_{\mathrm{pre}})^C)\\
\text{logarithmic complexity and slice cost}
  & (2+b_{\mathrm{pre}})^C\\
\text{positive fractions, scores, and errors}
  & \exp(-(2+b_{\mathrm{pre}})^C)\text{ from below}.
\end{array}
\]
We call a positive quantity with the last bound \emph{pre-large}.
None of these bounds may involve $\log L$, $\log R$, or the heights
of the $W_h$.  In particular, a pre-cost slice has polynomial
\emph{logarithmic} cost.

Finally choose the remaining parameters in the order
\begin{equation}\label{sampling:choice-order}
 \text{pre requirements}\ ;\quad
 \sigma\downarrow\ ;\quad A\uparrow\ ;\quad R\uparrow\ ;\quad
 \min_iH_i\uparrow.
\end{equation}
Here $0<\sigma<1$ is the perturbation parameter for the coefficient boxes
in the construction below.
The integer $L$ is selected using $W_h$ within
$A\le L\le A^{O_s(1+\sum_hd_h)}$.
The thresholds for $R$ and $H_i$ must work throughout that range, independently
of the subspace heights and of all polynomial coefficients.  Only a bounded
number of successive enlargements of the pre budgets is allowed.

\begin{proposition}[Constrained sampling]
\label{sampling:constrained}
Prescribe a finite collection of input caps, testing complexities, slice
bounds, and requested accuracies at the pre stage.
Assume the system has rank at least $R$ in Definition~\ref{sampling:rank},
with $R$ and the side lengths chosen sufficiently large in the order
\eqref{sampling:choice-order}.  For every tuple of centers, the construction
gives one ensemble satisfying all the prescribed requirements, with bounds
uniform in the centers, the inputs, and the path-dependent slice and test
choices within those requirements.  All later choices are made for this
prescribed collection; changing the requested accuracies permits a new
construction.  The ensemble consists of integer
affine maps $\psi(t)=x+Vt$ from a box
$\mathcal T=\prod_{j=1}^m\{0,\ldots,T_j-1\}$, where $m$ was fixed as a polynomial of $1+n+\sum_hd_h$
before all testing and accuracy budgets, $T_j\le L$, and some sides have
length $L$.  The ensemble has the following properties.
\begin{enumerate}
\item\label{sampling:property-geometry}
Each $\psi(\mathcal T)$ lies in $U$, and the functions $y_h\circ\psi$
are fully slow polynomial lifts while $\beta_h\circ\psi$ have integer
ordinary coefficients, all of degrees at most $h$.  Noninjective maps have
negligible probability.  If the centers are mixed independently with Haar
measure, the one-site law, with $t$ uniform, approximates the uniform law on
$U$ at the prescribed forward accuracy against every unit-bounded input.

\item\label{sampling:property-detection}
A permitted site twist is a bounded Lipschitz function of
$u/H,(y_h)_h$, multiplied by buffered chart cutoffs and by characters or
residue indicators of $(u,\beta_h)_h$ of a pre-bounded common modulus.
For fixed centers, suppose a pre-large fraction of sampled maps detects an
arbitrary pre-bounded function $\phi:U\to\C$, on a pre-cost slice of
$\mathcal T$, against a bounded local niltest of degree $r$, with
$0\le r\le s$, and pre
complexity.  The slices and niltests may be chosen separately for each map.
Then $\phi$ has a pre-large correlation on $U$ with a permitted twist times
a bounded ambient degree-$r$ niltest of pre complexity.

More precisely, write
\[
 X(\phi)=\E_\psi\sup_{B,F}
 \left|\E_{t\in B}\phi(\psi(t))F(t)\right|,
\]
where $B,F$ range over the prescribed slices and tests.  Every fixed
unit-bounded input admits an approximation in $X$, to any prescribed
pre accuracy, by a linear combination of the indicated ambient twisted
niltests.  The number of terms and the sum of the absolute coefficients
have pre bounds.

\item\label{sampling:property-transfer}
Let $T_*(u,\beta)$ be a fixed positive patch formula with $d_*$
further slots, in which $u$ has weight $1$ and $\beta_h$ has weight $h$,
and with pre-bounded complexity.  If $0\le f\le1$, $0\le\lambda\le1$,
and a pre-large fraction of maps have a pre-cost slice on which
\[
 \E_t(f(\psi(t))-\lambda)
 T_*(\psi(t),\beta(\psi(t)))\ge e^{-b_{\mathrm{pre}}},
\]
then, for fixed $\varepsilon>0$, a slice of $U$ supports a positive patch
of rank at most $d_*+\sum_h\dim W_h$ and an inverse-exponential-in-polynomial
score against $f-(1-\varepsilon)\lambda$.
The output complexity, slice cost, and logarithmic inverse score in this
item alone may charge the
logarithmic heights of bases of $W_h$ and $\log R,\log L$.
\end{enumerate}
There is also an unconstrained version of
items~\ref{sampling:property-geometry} and~\ref{sampling:property-detection}:
for every fixed
$\ell\ge\ell_0(s)$ one may prescribe the $\ell$ integer side lengths
in advance, provided they are sufficiently large and pairwise comparable.
In particular they may be prescribed distinct comparable primes.  Its
integer spatial law is dominated by a uniform product-box law at scales
$H_i$ and $H_i/L$ by a pre-bounded factor, as required for
Proposition~\ref{sampling:scalar-transfer}.  The choice of $\ell$ is fixed
before all testing and accuracy budgets; all constants may depend on it.
\end{proposition}

\begin{remark}[Proof and uniformity of the sampler]
Section~\ref{cube:sampler-proof} assembles all three properties for one
ensemble.
Corollary~\ref{sampling:normalization} supplies the geometric and one-site
conclusions in item~\ref{sampling:property-geometry}.
Lemmas~\ref{cube:pathwise-detection} and~\ref{cube:separation} give
detection and approximation in item~\ref{sampling:property-detection};
Proposition~\ref{cube:conditioned-transfer} gives
item~\ref{sampling:property-transfer}.  In
item~\ref{sampling:property-detection}, the output complexity is independent
of $\log L$, and the supremum permits tests and slices chosen separately
on each path.
\end{remark}

\subsection{The geometry of small lifts}

\begin{lemma}[Coordinates with bounds independent of subspace height]
\label{sampling:lattice-chart}
Let $W\subseteq\R^d$ be rational, let $\Lambda=W\cap\Z^d$, and let
$\Lambda'=\operatorname{proj}_{W^\perp}\Z^d$.  Choose a basis $v_i$ of
$\Lambda'$ with orthogonality defect and lengths at most $e^{O(d^2)}$,
and put $K_i=\max(1,\lceil |v_i|^{-1}\rceil)$.
In a small lift $y$ with $\operatorname{proj}_{W^\perp}y=\sum_iv_iz_i$,
the coordinates consisting of the orthogonal $W$ coordinates and the
numbers $z_i/K_i$ give forward and inverse linear maps of norm $e^{O(d^3)}$.
Moreover
\[
 e^{-O(d^3)}\le\frac{\covol(\Lambda)}{\prod_iK_i}
 \le e^{O(d^3)}.
\]
In these charts Haar measure on $W/\Lambda$ is Lebesgue measure in $W$
times counting measure in $z$, divided by $\covol(\Lambda)$.
\end{lemma}
\begin{proof}
The lattice $\Lambda'$ is generated by the projections of the standard
basis, all of length at most one; lattice reduction gives the displayed
bounds for a basis.  The vectors $K_iv_i$ have lengths between one and
$e^{O(d^2)}$, and the same orthogonality defect as the $v_i$.  Determinant
and singular-value bounds therefore give the claimed bounds for the linear
coordinate maps.

The subgroup $\Lambda$ is primitive in $\Z^d$, and the quotient
$\Z^d/\Lambda$ identifies with $\Lambda'$.  Lifting a basis of $\Lambda'$
and adjoining a basis of $\Lambda$ therefore gives a basis of $\Z^d$.
Its determinant, evaluated in the orthogonal splitting
$W\oplus W^\perp$, proves
$\covol(\Lambda)\covol(\Lambda')=1$.
The asserted covolume bound now follows from the determinant of the matrix
with columns $K_iv_i$.  Finally each local sheet is a translate of $W$;
translation invariance and the volume of a fundamental region of $\Lambda$
give the Haar formula.
\end{proof}

Choose integer lifts of the $v_i$ and let $l(z)$ be their integer linear
combination.  In Equation~\eqref{sampling:lift}, projection gives
$\operatorname{proj}\beta=-\sum_iv_iz_i$, so
$\beta+l(z)\in\Lambda$.  Let $w$ be its integer coordinates in a basis
of $\Lambda$.  The maps $\beta\mapsto(z,w)$ are integer linear maps.
Consequently conditions on $z,w$ modulo $M$ can be expressed by characters
whose coefficients are reduced modulo $M$, without paying for the heights
of the chosen lattice bases.  This observation does not bound their real
coordinate norms, and no such bound is needed for this modular statement.

\subsection{The coefficient tilt}

We now construct a product density on the polynomial coefficient arrays
and use it to weight the spatial affine maps.

Choose one $L$ within the allowed range so that, simultaneously in
all layers, no $K_i^{1/h}$ lies strictly between $L$ and $AL$. There are only
$\sum_hd_h$ relevant numbers, so a collection of more than that many
successive disjoint multiplicative gaps supplies such a choice, after
rounding endpoints with a fixed buffer.  Call a projection axis inactive if
$K_i\le L^h$, moderate if $L^h<K_i\le L^{C_s}$, and enormous otherwise.
An \emph{active} integer axis means a moderate or enormous projection axis.

Give each continuous or projection axis a number $J_0$, fixed as a
polynomial of $1+n+\sum_hd_h$ before all accuracy budgets, of separate blocks of $h$ parameter columns.  Use side length $L$ in these
blocks, except for an inactive axis, where it is
$\max(1,\lfloor cK_i^{1/h}\rfloor)$.  Include further disjoint kernel
blocks, each with $n$ columns.  In normalized spatial coordinates the kernel
matrices lie in small neighborhoods of fixed invertible matrices, with
pre-bounded inverse norms.  The number of kernel blocks is likewise
structural.  Choose these counts to dominate the structural dimension
multiples in the modular-rank estimate.  Choose $J_0$ large enough for
Lemma~\ref{cube:product-block}; its proof gives a threshold polynomial
in the structural dimensions alone.  No requested precision is allowed
to increase these block counts or $m$; precisions are handled afterwards
by choosing widths, perturbation size, and the later scales.

Subtract the chosen representative of each center $c_h$ only from the
constant coefficient.  The polynomial lift has an
independent small constant shift on each normalized axis, with widths
chosen at the pre stage and independent of $\sigma,L$.  Its main
nonconstant part is a sum of the products of the $h$
parameters in the dedicated blocks.  The coefficient is $1$ on an inactive
integer axis (or $0$ for a bounded axis), an integer in a small interval
of length comparable to $K_i/L^h$ on an active integer axis, and a real
number in an interval of length comparable to $L^{-h}$ on a continuous
axis.  All remaining nonconstant coefficients vanish on inactive and
moderate axes.  On continuous and enormous axes give these remaining
coefficients smooth perturbations of normalized
width comparable to $\sigma/T^\nu$.  Constants and widths are chosen,
after the dimensions, so the sum of all monomial contributions stays well
inside the prescribed chart throughout the parameter box.

All principal coefficients in different dedicated blocks are independent.
On continuous and active integer axes choose their normalized supports
inside $[\gamma_h,2\gamma_h]$, with $\gamma_h>0$
fixed after the structural block counts and chart widths.  The normalization
is multiplication by $L^h$ on a continuous axis and by $L^h/K_i$
on an active integer axis.  The constants $\gamma_h$ are small enough for the
preceding chart-containment requirement.  Thus even after these
coefficients are fixed, their normalized magnitudes are bounded below
by a pre-large constant.  Continuous choices have fixed smooth
buffered densities after this normalization.

For an integer choice with center $a$ and scale $S$, use the probability
mass
\[
 p(k)=\frac{\psi((k-a)/S)}{\sum_{j\in\Z}\psi((j-a)/S)},
\]
where $\psi$ is a fixed nonnegative, nonzero smooth bump supported in
$(-1,1)$, and $a,S$ are chosen so its rescaled support lies strictly inside
the selected interval.  This convention also applies to the small
perturbation coefficients on enormous axes.  When $S=K_i\delta$ and
$K_i\delta$ is sufficiently large, the denominator is comparable to
$K_i\delta$, and
\[
 q(u)=\frac{K_i\psi((K_i u-a)/(K_i\delta))}
                 {\sum_{j\in\Z}\psi((j-a)/(K_i\delta))}
\]
interpolates $K_i p(k)$ at $u=k/K_i$, with derivatives of order $r$
bounded by $C_r\delta^{-r-1}$ independently of $K_i$.
Every enormous-axis scale satisfies this lower bound by the later
choice of $L$.  On inactive or moderate axes the permitted
singleton interpolation instead costs $O(\log K_i)=O_s(\log L)$.
For constant projection coordinates use this law centered at zero at a
fixed pre normalized width whenever $K_i$ times that width is sufficiently
large; otherwise choose the constant zero.  In the latter case $K_i$ is
pre-bounded.  A pre-bounded inactive axis may likewise have all its
nonconstant coefficients zero.  In particular when $W_h=\{0\}$ every
coefficient is zero.  Constant continuous coordinates have independent
smooth densities of fixed pre width about zero.  These constant choices
are independent of all nonconstant coefficients.
All remaining continuous perturbations are sampled as their center
plus their stated width times an independent variable with a fixed
smooth buffered density.  This specifies the probability laws, not
only their supports.

To verify normalization, for one coefficient write $u=y_W$ and let
$f(u)$ be its normalized continuous density and $p_i(z_i)$ its independent
projection masses.  Lemma~\ref{sampling:lattice-chart} gives its density
relative to Haar explicitly as
\[
 \covol(\Lambda_h)f(u)\prod_i p_i(z_i),
 \qquad
 d\mathrm{Haar}=\frac{du\,d\#z}{\covol(\Lambda_h)}.
\]
Its integral is exactly one.  For a constant coefficient, write this density as
\[
 \frac{\covol(\Lambda_h)}{\prod_iK_i}\,f(u)\prod_i\bigl(K_ip_i(z_i)\bigr).
\]
Every factor has a pre cap: the smooth constant widths are pre, and a
singleton constant was used only when $K_i$ is pre-bounded.  Thus the
constant-coefficient density has a pre cap independently of subspace height.
The supports were chosen so their full
Cartesian product lies inside the buffered chart, so no truncation or
further normalization is required.  Singleton projection masses select
discrete local sheets and do not destroy absolute continuity.  Multiplying
these independent laws over coefficients defines a density $\mathcal D$
of integral one on the
product coefficient torus
\[
 \prod_{h=1}^s\prod_{|\nu|\le h}W_h/\Lambda_h.
\]
Let $\rho(x,V)$ be a smooth product law on integer choices, with $x$ nearly
uniform on an interior box.  Its kernel-block matrices have the
pre-bounded neighborhoods and inverse bounds specified above.  Every
nonkernel entry is instead supported in
\begin{equation}\label{sampling:nonkernel-support}
 |V_{ij}|\le \frac{c\sigma H_i}{mT_j},
\end{equation}
with a fixed smooth profile after rescaling by this width, independently
of $x$, the kernel entries, and all other nonkernel entries.  Consequently
the total normalized displacement from all nonkernel columns is
$O(\sigma)$ on the parameter box.  The laws of $x$ and the kernel
entries do not depend on $\sigma$.  Normalized derivative and cap
logarithms for the whole spatial law are polynomial in the pre budget
and $\log\sigma^{-1}$; these are allowed later costs.  The actual
sampling law is
\begin{equation}\label{sampling:tilted-law}
 \frac1Z\rho(x,V)\,
 \mathcal D\bigl(\operatorname{coeff}(C_h(x+Vt))_{h}\bigr),
 \qquad
 Z=\E_\rho\mathcal D\bigl(\operatorname{coeff}(C_h(x+Vt))_{h}\bigr).
\end{equation}
For the unconstrained version set $\mathcal D=1$, omit all lift and
dedicated-block conditions, and use any fixed number
$\ell\ge\ell_0(s)$ of columns.  Choose independent smooth integer
entries of $V$ in small but nondegenerate intervals at scales $H_i/T_j$.
Their normalized widths are chosen after the testing budgets and have
pre-bounded inverse logarithms.  Together with the similarly chosen $x$
law, this is dominated by a uniform product-box law by a pre-bounded
factor.  The bounds use only comparability of the prescribed $T_j$ to
$L$, not equality or compositeness; hence prescribed distinct prime sides
are allowed.  All columns serve as the free columns in the subsequent
cube argument.  Taking $\ell$ above its fixed degree-dependent threshold
makes the exceptional dependent parameter cubes negligible, as in that
argument.  No coefficient tilt or special axis lengths need to be selected
in this version.

The identity $\int\mathcal D=1$ on the abstract coefficient torus does
\emph{not} imply $Z=1$; the latter requires equidistribution of the integer
spatial choices.  On the support of the tilt, subtracting $c_h$ from the
constant coefficient and then the prescribed small coefficient lifts
leaves integer ordinary coefficients.  Their polynomial sum defines
$\beta_h\circ\psi$, and the sum of the small lifted monomials defines
$y_h\circ\psi$.  The latter stays inside the uniqueness chart at every
parameter site and has every coefficient bounded at its actual $T^\nu$
scale.  Thus, whenever $Z>0$, these polynomials agree with the sitewise
lifts in Equation~\eqref{sampling:lift} and give the asserted fully slow
and integer-coefficient identities exactly.

Lemma~\ref{sampling:lattice-chart} gives the bound for ambient
Fourier approximations.  After multiplying an integer-axis probability mass
by $K_i$, its normalized density has height-independent bounds.  A singleton
restriction on an inactive or moderate axis can be extended by a smooth
window of width less than $1/(2K_i)$, with logarithmic cost $O(\log L)$.
On an enormous axis every prescribed interval contains many grid points
once $C_s>s$ and $L$ is sufficiently large in terms of $\sigma^{-1}$.
Its smooth interpolation then has costs independent of $K_i$.  Smooth
windows supported strictly within the ambient charts can be periodized and
approximated by product Fourier kernels.  This gives costs polynomial in
pre budgets, $\log\sigma^{-1}$, $\log L$, and inverse error logs.
The pre bounds in item~\ref{sampling:property-detection} of
Proposition~\ref{sampling:constrained} follow from the cube-image comparison
in Proposition~\ref{cube:comparison} and the approximation argument in
Lemma~\ref{cube:separation}.

\subsection{Fourier modes and Boolean cube evaluation}

\begin{lemma}[The ideal of an affine Boolean cube]
\label{sampling:cube-ideal}
Let $t_\omega=t_0+\sum_{j=1}^q\omega_je_j\in\Z^m$, where the
$e_j$ are linearly independent, and put $T_\omega=(1,t_\omega)$.
For $h\ge1$, the vector space of homogeneous degree-$h$ real polynomials
vanishing at every $T_\omega$ is spanned by products of $h$ rational linear
forms such that, at each site $T_\omega$, at least one factor vanishes.
The vanishing factor may depend on $\omega$.
If $q,h$ are bounded and $|t_\omega|\le L$, the forms can be chosen with
heights bounded by $(2+m+L)^{O_{q,h}(1)}$.
\end{lemma}
\begin{proof}
A rational change of coordinates sends the points to
$(a_0,a_1,\ldots,a_q,z)=(1,\omega_1,\ldots,\omega_q,0)$.
One can obtain this change using a nonzero $q\times q$ minor of the matrix
of the $e_j$; Cramer's rule gives the asserted height bound.

Reduce a homogeneous monomial by the relations $z_i=0$ and
$a_i^2=a_0a_i$.  Its remainder is a linear combination of
$a_0^{h-|S|}\prod_{i\in S}a_i$, where $|S|\le h$.
The functions $\prod_{i\in S}\omega_i$ on $\{0,1\}^q$ are linearly
independent, by successive evaluation on indicator vectors of subsets.
Therefore a polynomial vanishing on the cube has zero remainder, and lies
in the degree-$h$ part of the ideal generated by
\[
 z_i,\qquad a_i(a_i-a_0)\quad(1\le i\le q).
\]
Multiplying each generator by monomials gives the required spanning
products.  The height bound is preserved under this bounded-degree
construction and the inverse coordinate change.
\end{proof}

\begin{lemma}[Algebraic reduction of non-site modes]
\label{sampling:mode-removal}
Let $W\subseteq\R^d$ and let $\lambda$ be a row functional on the
coefficients of a homogeneous $W$-valued degree-$h$ polynomial in $T$.
If $\lambda$ does not factor through evaluation at the points in
Lemma~\ref{sampling:cube-ideal}, there are rational linear forms
$a_i\cdot T$, $1\le i\le h$, whose product vanishes at every site, such
that
\[
 \lambda'(w)=\lambda\left(w\prod_{i=1}^h(a_i\cdot T)\right)
\]
is a nonzero functional on $W$.  Its coefficient bounds depend only on
those of $\lambda$ and the parameter tuples, and not on the height of $W$.
\end{lemma}
\begin{proof}
A functional factors through a linear map if and only if it vanishes on
its kernel.  The kernel of site evaluation on $W$-valued polynomials is
$W$ tensored with the scalar kernel described in
Lemma~\ref{sampling:cube-ideal}.  If $\lambda$ does not vanish there,
it is nonzero on $wP$ for one of the spanning products $P$ and some $w\in W$.
This gives the assertion.  The construction of $P$ involved only the
parameter tuples, which proves the last statement.
\end{proof}

\begin{lemma}[Quantitative removal of non-site modes]
\label{sampling:analytic-removal}
Assume Lemma~\ref{prelim:step-drop} and that the polynomial system has
rank at least $R$ in Definition~\ref{sampling:rank}.  Prescribe a Fourier coefficient and
frequency budget, a common cover modulus, spatial smoothness bounds, and
an inverse-error log.  Include $m,n,\log L$ in a budget $b$ containing
these inputs; for Equation~\eqref{sampling:nonkernel-support} include
$\log\sigma^{-1}$ as well.  The spatial law may also be restricted to
bounded-modulus
congruences.  For a fixed cube of bounded order with independent parameter
differences, insert arbitrary unit-bounded functions of its sites in a
spatial expectation.  Every coefficient Fourier mode that fails to factor
through real site evaluation on at least one $W_h$ has expectation at most
the prescribed error, provided $R\ge\mathcal B(b)$ and the parent sides
are sufficiently large $\mathcal B(b)$.  Neither bound charges the heights
of $W_h$.  The same conclusion holds for one site.
\end{lemma}
\begin{proof}
It suffices to treat one mode and one spatial congruence; the specified
Fourier coefficient mass absorbs their sums.  Choose the highest layer
$h$ at which the mode does not factor.  Every higher-layer phase is a real
linear combination of the values $C_j(vT_\omega)$ and can therefore be
absorbed into the arbitrary bounded site functions.  Such real linear
coefficients can be chosen by a rational right inverse of the parameter
evaluation matrix, with the prescribed height bounds: the kernel condition
holds on $W_j$, so composing with that right inverse agrees with the mode
there.  No basis of $W_j$ is needed.

By Lemma~\ref{sampling:mode-removal}, choose rational forms $a_i\cdot T$
whose product vanishes at every site and whose contraction gives a row
$\lambda'$ nonzero on $W_h$.  Clear their denominators.  Replacing the
forms by integer multiples changes $\lambda'$ by a bounded nonzero rational
factor; all numerator and denominator logs remain polynomial in $b$.
Write the remaining phase as $P(v)$, of degree at most $h$ in $v=(x,V)$.

Take independent uniform integer vector shifts $w_i$ in boxes of side
lengths $\widetilde H_a$, with a common stride divisible by the spatial
congruence modulus.  Choose $\widetilde H_a$ to be the original $H_a$
divided by a sufficiently large $\mathcal B(b)$ factor.  This factor
also dominates the reciprocal normalized widths, all $T_j\le L$, the
integer forms, and the congruence stride.  Every displacement
$\sum_iw_ia_i^{\mathsf T}$ is therefore small compared with each actual
spatial width, including $c\sigma H_a/(mT_j)$ on nonkernel columns.
The total variation distance between the smooth spatial law
and each such translate is smaller than the prescribed error.  This
follows by integrating its normalized derivative bound; integer rounding
and the fixed stride cause no loss when the sides are sufficiently large.
Thus if the original mode expectation has modulus at least $\delta$,
averaging these translations changes it by at most $\delta/2$.
The triangle inequality then supplies one base tuple $v$ for which
\[
 \left|\E_{w_1,\ldots,w_h}
 e\left(P\left(v+\sum_iw_ia_i^{\mathsf T}\right)\right)
 \prod_\omega F_\omega\left(
    vT_\omega+\sum_iw_i(a_i\cdot T_\omega)\right)\right|
 \ge\delta/2.
\]
Extend the bounded site functions outside $U$ by zero.  Each site factor
misses at least one $w_i$, because $a_i\cdot T_\omega=0$ for at least
one $i$; assign it to one such variable.  Applying
Cauchy--Schwarz once
in each independent variable gives an absolute phase average at least
$(\delta/2)^{2^h}$ on independent pairs $w_i,w'_i$.
All terms of degree less than $h$ vanish.  The degree-$h$ term is precisely
\[
 h!\,\operatorname{Pol}(\lambda'C_h^{\mathrm{top}})
 (w_1-w'_1,\ldots,w_h-w'_h),
\]
up to sign.  This identity follows by taking the coefficient of the product
of the $h$ shift scalars in
$C_h^{\mathrm{top}}(\sum_iw_i(a_i\cdot T))$ and applying the mode to its
$T$ coefficients.  In particular the resulting phase is independent of
the chosen base tuple $v$.

Apply Lemma~\ref{prelim:step-drop} to this polynomial on the last-layer
torus: take $G=\R$, $\Gamma=\Z$, and $G_1=\cdots=G_h=\R$,
$G_{h+1}=0$, with character $e(u)$ of frequency one.  The associated
graded algebra has only grade $h$, so the fast algebra killed by that
character is zero.  There are $2hn$ integer variables, its inverse bias
log is $O_h(1+\log\delta^{-1})$, and all individual sides are sufficiently
large after the preceding shrinkage and stride.  Thus the exact hypotheses
of the lemma hold without a bound on the phase coefficients.  Hence
the homogeneous polynomial just displayed is a sum of a slow polynomial
and a rational polynomial with common denominator $\mathcal B(b)$.
Set $w'_i=0$ in the homogeneous identity and identify all $w_i=w$.
There are only polynomially many coefficients, so this restriction preserves
slow bounds at the $\widetilde H_a$ scales, up to $\mathcal B(b)$.
The diagonal left side is $h!\lambda'C_h^{\mathrm{top}}(w)$.
Clear the denominator of $\lambda'$ and the shift strides, absorbing
$h!$ into the integer row and any resulting rational denominator, and
restore the original $H_a$ scales.  We obtain an integer row $a$ of norm
$\mathcal B(b)$, nonzero on $W_h$, for which $(aC_h)^{\mathrm{top}}$ is
slow plus rational with these same permitted bounds.  Choosing $R$ larger
than all these bounds contradicts Definition~\ref{sampling:rank}.
This proves the removal estimate uniformly for each mode before its
Fourier coefficient is summed.
\end{proof}

\begin{remark}[Real factorization and torus covers]
Factoring through evaluation as a real linear map is not necessarily the
same as factoring through evaluation on a torus.  For example, the character
$x\mapsto e(x)$ does not factor through $x\mapsto2x$ on $\R/\Z$.
For rational evaluation matrices one may clear the denominators of a
rational right inverse by a finite common cover.  Surviving real-factoring
modes must be read on that cover, and the coefficient density must be lifted
to the same cover.  The required denominators depend on the parameter
matrix, not on $W$.  Ignoring this distinction would invalidate the mode
removal argument.
\end{remark}

\begin{corollary}[Coefficient equidistribution and one-site consequences]
\label{sampling:normalization}
Assume the hypotheses and symbol estimate of
Lemma~\ref{sampling:analytic-removal}.  The coefficient arrays of
$C_h(x+Vt)$ are equidistributed to any prescribed accuracy against the
bounded-cost Fourier approximations used in the tilt, jointly with the
spatial smooth weights and bounded-modulus congruences.  Consequently
$Z=1+o(1)$ in Equation~\eqref{sampling:tilted-law}.  Mixed centers give the
one-site approximation in item~1 of
Proposition~\ref{sampling:constrained}; with fixed centers the one-site
law has negligible excess mass above a pre-bounded density relative to
uniform measure on $U$.  Noninjectivity has negligible probability.
\end{corollary}
\begin{proof}
Without site insertions, the proof of
Lemma~\ref{sampling:analytic-removal} removes every nontrivial coefficient
mode: choose the highest nonzero layer and a product of coordinate forms
on which its coefficient functional is nonzero.  The vanishing-at-sites
condition is unnecessary because there are no site factors to eliminate.
Only the trivial characters on each $W_h/\Lambda_h$ remain, exactly as in
product Haar integration.  This proves the stated equidistribution and,
using the Fourier approximation of $\mathcal D$, the normalization claim.
This estimate is uniform in the centers: translating a constant-coefficient
law changes Fourier coefficients only by phases, preserving all budgets.
It therefore also permits averaging the normalized laws $\rho\mathcal D/Z$
over those centers.
The same reasoning works on a common integer cover.  It also proves Haar
equidistribution of the coefficient arrays obtained by expanding
$C_h(x+\sum_{i=1}^q\omega_i d_i)$ on $\omega\in\{0,1\}^q$ and
reducing $\omega_i^2=\omega_i$.  These squarefree coefficient arrays
are the \emph{Boolean jets}; their dimensions are recorded in
Equation~\eqref{cube:jet-dimension}.  In ambient averages over the base
$x$ and independent differences $d_i$, the preceding highest-layer
argument applies to every nonzero character of these arrays.
Polytope restrictions
that the cube lie in $U$ can be approximated by spatial smooth cutoffs,
with boundary errors bounded by elementary box counting.

For one site, Lemma~\ref{sampling:analytic-removal} keeps only modes
factoring through that evaluation.  Averaging the centers removes every
nonzero frequency on the constant coefficient of each layer.  A mode
factoring through one evaluation has its site row also as its constant
coefficient row, so the only surviving mode is trivial.  The unweighted
one-site spatial law approximates uniform measure on $U$: choose $x$
nearly uniform on a box omitting a prescribed small boundary fraction and
choose $V$ so that the maximal site displacement is a still smaller fraction
of each $H_i$.  The corresponding translation and boundary errors are
arbitrarily small at the forward pre accuracy.  This proves the mixed-center
assertion against arbitrary bounded site functions.

With fixed centers, the surviving weight is the conditional Haar average
of the coefficient density given the site value.  Its constant coefficients
are independent of the other coefficients and have a pre-bounded density
relative to Haar, by the fixed pre widths of the constant laws and
Lemma~\ref{sampling:lattice-chart}; convolution with
the remaining coefficients cannot increase that cap.  The unweighted
spatial pushforward has a pre-bounded density because of the full-scale
$x$ shift, even after conditioning on all of $V$.  In particular this cap
does not use the $\sigma$-dependent joint spatial density or the late
supremum bound of $\mathcal D$.  Thus for every measurable set $E\subseteq U$ the one-site law
obeys $\Pr(E)\le C\Pr_U(E)+\eta$, where $C$ is pre-bounded and $\eta$
is arbitrarily small in the later choice order.  Apply this to the set
where its density exceeds $2C$ to bound the total excess mass there by
$2\eta$.

Finally, for a fixed nonzero difference $a=t-t'$ a collision requires
$Va=0$.  In one row, after all but a column with $a_j\ne0$ are fixed,
at most one value of $V_{ij}$ is possible.  Its probability is at most
$\mathcal B(b)L/(\sigma H_i)$ under the smooth spatial law, allowing a
nonkernel pivot in Equation~\eqref{sampling:nonkernel-support}.  A union bound over
at most $L^{2m}$ pairs, followed by the available supremum bound on the
tilt and $Z=1+o(1)$, makes the collision probability arbitrarily small
when the parent sides are sufficiently large.  All losses have polynomial
logarithms in the already chosen budgets.
\end{proof}

\begin{corollary}[Finite residues under the actual tilted law]
\label{sampling:finite-residues}
Fix centers and an error $\eta>0$.  Prescribe a finite family of joint
residue vectors, each consisting of polynomially many spatial residues
and integer coefficient lifts $\beta_h$ or their integer-linear
coordinates $z,w$, modulo a common modulus $M$ for that vector.
The maximum of $\log M$ and the logarithm of the family size have prescribed
bounds.  Include these bounds in the Fourier budget of
Lemma~\ref{sampling:analytic-removal}.
For every vector in the family, its joint distribution under
Equation~\eqref{sampling:tilted-law} is within $\eta$ in total variation
of the following comparison law: sample the coefficient-chart residues
on the Haar cover with density $\mathcal D$, and independently sample
the corresponding spatial residues under the unweighted law $\rho$.
Exact independence in this comparison law is not an assertion of exact
independence under the tilted law.

This permits a family of separate tests modulo prime powers at most $Q$
and modulo prescribed bounded products of witness prime powers.
It does not assert a joint approximation modulo the product of all primes
at most $Q$, whose logarithm need not have the permitted bound.
The rank threshold and parent side cutoff may depend on the maximum
individual $\log M$, the log family count, and $\log\eta^{-1}$;
previously fixed bounds for the chart maps and coefficient-law parameters
do not.
\end{corollary}

In particular this is a joint finite-residue comparison.  It does not
assert conditional uniformity under the actual law after fixing an exact
real lift value.  A bounded insertion depending on the site is handled
by Lemma~\ref{sampling:analytic-removal} with that insertion retained;
the finite-residue conclusion alone does not remove it.

\begin{proof}
Treat one joint residue vector at a time.  Lift every coefficient torus
to an integer cover divisible by its modulus $M$ and by any evaluation
denominators.  Under Haar measure
on $W_h/M\Lambda_h$, the deck coordinate in $\Lambda_h/M\Lambda_h$
is uniform and independent of its image in $W_h/\Lambda_h$.  In the chart
coordinates this says that the $w_{h,\nu}\bmod M$ coordinates are
independent uniform variables across coefficients, while the $z_{h,\nu}$
retain their prescribed projection-grid laws under the tilt.  Fixing centers
and any prescribed finite collection of chart pieces depending only on
the small lifts and projection coordinates does not change this
deck-coordinate statement.  The Fourier approximation below concerns
$\mathcal D$ times residue indicators; any further continuous chart
restriction is taken with a smooth buffered cutoff.  Within a buffered
coefficient chart, the integer lift residue is locally constant.
The functions $\mathcal D$ times its residue indicators consequently
admit ambient smooth Fourier approximations on the cover with the same
permitted costs, enlarged polynomially in the logarithm of its modulus.
For $z,w$ this uses that they are integer-linear functions of $\beta$;
the character coefficients can be reduced modulo the common modulus.
There is no dependence on the heights of the lattice bases.  Here is an
explicit ambient extension.  The map $W/M\Lambda\to\R^d/M\Z^d$ is
injective, since $W\cap M\Z^d=M\Lambda$.  Around each
$c+\beta$, $\beta\in(\Z/M\Z)^d$, use the disjoint buffered chart
$c+\beta+y$, $|y|_\infty<1/4$.  On that chart the density is a product
of smooth functions of $\operatorname{proj}_Wy$ and the normalized
projection coordinates $z_i/K_i$, multiplied by
$\covol(\Lambda)/\prod_iK_i$.  Lemma~\ref{sampling:lattice-chart}
bounds both these linear maps and the last factor independently of the
height.  Interpolate discrete masses after multiplying them by $K_i$: a
singleton window costs at most a power of $L$ on a moderate or inactive
axis, while an enormous axis has many grid points in every interval, so
its interpolation costs only inverse normalized interval widths.  Multiply
the whole chart by the required residue mask of its index $\beta$.
Disjoint buffered supports preserve the supremum and Lipschitz bounds
when these functions are periodized on $\R^d/M\Z^d$.  The same
construction in the product of coefficient charts has polynomial
dimension.  Fourier smoothing on this ambient torus therefore gives
uniform approximation, with frequency and coefficient-mass logarithms
polynomial in the permitted budgets and $\log M$, independently of
subspace heights.

Apply coefficient equidistribution with each imposed spatial residue class.
The shortest spatial entry width is at least
$\sigma H_i/\mathcal B(b)L$.  Choosing the parent sides after the
individual modulus $M$ makes counting in its residue classes arbitrarily
accurate, also for these nonkernel entries.  This uses only the permitted
late inverse widths and $\log M$, and preserves the product spatial law.
It compares the unnormalized expectation of every joint residue indicator
to its product Haar-cover integral times the unweighted spatial probability,
with an arbitrarily prescribed absolute error.  The number of atoms of
the finite joint residue space has bounded logarithm: it is at most the
polynomial number of coordinates times $\log M$.  Choose the error per atom smaller than
$\eta$ divided by this number of atoms and by a fixed constant.
Summing absolute errors gives total variation at most $\eta/2$ for the
unnormalized laws.  The normalization $Z=1+o(1)$ contributes at most the
remaining $\eta/2$.  This proves the assertion.  In particular it compares
the actual integer-spatial model with the abstract tilted model.  The
estimates are uniform over the prescribed family by using the maximum
individual budget; to sum errors over the family, first divide the desired
error by its size.  One never forms the possibly enormous common modulus
of all family members.  Haar approximation on continuous variables alone
would not imply the stated finite-residue conclusion without this argument.
\end{proof}

More explicitly, if $b_{\mathrm{pre}}$ is the prescribed initial budget
and $\eta$ the desired comparison error, the Fourier approximation budget
can be taken as
\[
 \begin{aligned}
 b_0=\mathcal P\bigl(&b_{\mathrm{pre}}+\log\sigma^{-1}
             +\log L+\log\eta^{-1}\\
             &+\max\log M+\log(2+\text{family size})\bigr).
 \end{aligned}
\]
The rank cutoff is $R_*=\exp(\mathcal P(b_0))$.  Replace $\log L$ here
by its upper bound $O_s(1+\sum_hd_h)\log A$ to choose $R_*$ uniformly
before the subspaces select $L$.

The Fourier and equidistribution costs in this subsection are allowed to
use $\log L$ and $\log\sigma^{-1}$.  Section~\ref{cube:section} supplies the
pre bounds for the pathwise testing seminorm in
item~\ref{sampling:property-detection} of
Proposition~\ref{sampling:constrained}, through
Proposition~\ref{cube:comparison} and Lemma~\ref{cube:separation}.
Section~\ref{cube:transfer-section} supplies fixed-patch transfer in
item~\ref{sampling:property-transfer}, through
Proposition~\ref{cube:conditioned-transfer}.
